\paragraph{Missing step.} The full target of Problem~C is the full quadratic scalar operator on the locked vacuum in a constructed frame: the canonically normalised kinetic coefficient of the pure-trace scalar, including the induced Einstein--Hilbert contribution, its mixing with \(T_{\mu\nu}\), its non-minimal curvature coupling, its matter coupling, and its pole mass, or a ratio \(m_\psi/v\), as a derived output, together with the coupling that can be compared with Cassini bounds, and the phenomenological comparison of Chapter~\ref{ch:agenda}, defining and testing either a sufficiently heavy \(m_\psi\) or the scalar-tensor outcome, with Assumption~\ref{ass:gap-psi} discharged. On the bare data of Presentation~A this operator is not yet definable: three of the inputs it needs are not derived. (i)~A frame. The transformation law of Equation~\eqref{eq:Phi-transform} is a field-content choice of Presentation~B and is not derived from Presentation~A alone; the stated Repair~I map to an orthonormal frame is obstructed above; the remark following Definition~\ref{def:vierbein} records that a corrected Repair~I would need additional background tensors or an independent coframe, which is part of Problem~C; and the route actually used posits the frame through Assumptions~\ref{ass:nondeg} and~\ref{ass:frame-lock}. (ii)~\(G_{\mathrm{eff}}\). The coefficient \(Z\) of Equation~\eqref{eq:Z-sigma} is written in terms of \(G_{\mathrm{eff}}\), whose derivation is Problem~F, so the scalar wave-function renormalisation ties Problem~C to Problem~F. (iii)~The constraint structure of the massless sector. The trace-constrained Fierz--Pauli structure of the quadratic tensor block is not shown to follow from the NAMCG quadratic block, and whether the \(\psi\) sector of Chapter~\ref{ch:massless-tensor} supplies it is unresolved and is a dependency of Problem~C; the \(\psi\) of Chapter~\ref{ch:massless-tensor} is the trace of the metric fluctuation \(h_{\mu\nu}\), distinct from the amplitude \(\psi\) of the modes \(\psi\Lambda_A\) of \(K\) in the Problem~F entry. With these three inputs supplied, the recorded results of this section (the tree radial mass of Lemma~\ref{lem:msigma}, the wave-function factor \(Z\) of Equation~\eqref{eq:Z-sigma} at the stated truncation, and the off-shell trace--tensor mixing term) are computations within posited frame data; the loop-level kinetic coefficient, the mixing, the non-minimal coupling, the matter coupling and the pole mass are not outputs of the matrix theory, and in the conjugate-coordinate class the \(q=0\) non-derivative sector of the static scalar channel carries a bounded obstruction recorded in Section~\ref{sec:status-F}.